% ==================================================================
% id:       func-002
% topic:    מספרים ממשיים
% subtopic: אינדוקציה מתמטית
% solution: full
% ==================================================================
\textbf{הוכיחו באינדוקציה את הטענות הבאות:}
\begin{enumerate}
\item
לכל
$x \in (-1, \infty)$
ולכל $n$ טבעי מתקיים
$(1 + x)^n \geq 1+nx$.
\begin{solbox}
נוכיח באינדוקציה כי לכל
$x \in (-1, \infty)$
ולכל $n$ טבעי מתקיים
$(1 + x)^n \geq 1+nx$.

נקבע $x$ בתחום.

בסיס האינדוקציה: נבדוק ל
$n=1$:
ואכן
$(1 + x)^1 = 1 + x  = 1+ 1 \cdot x$.

שלב האינדוקציה: נניח ל-$n$ ונראה ל
$n + 1$:

$$(1 + x)^{n + 1} = (1 + x)^n \cdot (1 + x)  \geq $$
נשתמש בהנחת האינדוקציה ובכך שעל פי ההנחה,
$1 + x$
חיובי, ולכן ניתן לכפול בו אי שיוויון
$$\geq (1+nx) \cdot (1 + x) = 1 + (n+1)x + nx^2 \geq 1 + (n+1)x$$
המעבר האחרון נכון מכיוון ש
$nx^2 \geq 0 $
לכל $n$ טבעי ולכל $x$ ממשי.

הוכחנו את שלב האינדוקציה, ולכן הטענה נכונה לכל $n$ טבעי.
\end{solbox}

\item $2^n \geq n^3$
לכל
$n \geq 10$
טבעי.
\hint{בשלב המעבר השוו את $\frac{2^{n+1}}{(n+1)^3}$ ל-$\frac{2^n}{n^3}$, והשתמשו ב-$n\geq10$ כדי לחסום את היחס $\frac{n^3}{(n+1)^3}$ מלמטה.}
\begin{solbox}
נוכיח באינדוקציה כי
$2^n\geq n^3$
לכל
$n\geq10$
טבעי.
\begin{itemize}
\item
שלב הבסיס: נציב
$n=10$
ונקבל
$$,2^{10}=1024>1000=10^3$$
כלומר, הטענה נכונה עבור
$n=10$.
\item
שלב ההנחה: נניח כי אי-השוויון הנתון נכון לאיזשהו
$n\geq10$,
כלומר,
$2^n\geq n^3$
עבור איזשהו
$n\geq10$.
לצורך ההוכחה שנציג נוח יותר לרשום
$2^n\geq n^3$
בצורה שקולה
$\frac{2^n}{n^3}\geq1$.
\item
שלב היסק: מתקיים
\begin{align*}
\frac{2^{n+1}}{(n+1)^3}&=2\cdot\frac{2^n}{(n+1)^3}=2\cdot\frac{2^n}{{\color{red}n^3}}\cdot\frac{{\color{red}n^3}}{(n+1)^3}\\
&\underset{\text{הנחת האינדוקציה}}{\geq}2\cdot1\cdot\frac{n^3}{(n+1)^3}=2\cdot1\cdot\lr{1-\frac{1}{n+1}}^3\\
&\geq 2\cdot\lr{1-\frac{1}{10+1}}^3=2\cdot\frac{10^3}{11^3}=\frac{2000}{1331}>1,
\end{align*}
כנדרש.
\end{itemize}
\end{solbox}
\item $n! < \frac{(n+1)^n}{2}$
לכל
$n > 1$
טבעי.
\hint{בשלב המעבר כתבו $(n+1)!=n!\cdot(n+1)$ והשתמשו בהנחת האינדוקציה.}
\begin{solbox}
נראה שמתקיים
$n!<\frac{(n+1)^n}{2}$
לכל
$n>1$
טבעי.
\begin{itemize}
\item
בסיס האינדוקציה: מתקיים
$$,2!<\frac92=\frac{(2+1)^2}{2}$$
כלומר, האי-שוויון תקף עבור
$n=2$.
\item
שלב ההנחה: נניח כי
$n!<\frac{(n+1)^n}{2}$
עבור איזשהו
$n\geq2$.
\item
שלב היסק: מתקיים
\begin{align*}
(n+1)!&=n!\cdot(n+1)\\
&\underset{\text{הנחת האינדוקציה}}{<}\frac{(n+1)^n}{2}\cdot(n+1)\\
&=\frac{(n+1)^{n+1}}{2}<\frac{(n+2)^{n+1}}{2},
\end{align*}
כנדרש. המעבר האחרון נכון כי הגדלנו את הבסיס של החזקה מ-$n+1$ ל-$n+2$
(שניהם חיוביים) ולא שינינו את המעריך $n+1$.
\end{itemize}
\end{solbox}
\item
לכל טבעי
$n$:
$\ds 1^2+2^2+\cdots+n^2=\frac{n(n+1)(2n+1)}{6}$
\begin{solbox}
נראה באינדוקציה כי לכל
$n$
טבעי מתקיים
$$.1^2+2^2+\cdots+n^2=\frac{n(n+1)(2n+1)}{6}$$
\begin{itemize}
\item
בסיס האינדוקציה: נציב
$n=1$
ונקבל
$1^2=1=\frac{1\cdot(1+1)(2\cdot1+1)}{6}$,
כלומר, הטענה נכונה עבור
$n=1$.
\item
שלב ההנחה: נניח כי
$1^2+2^2+\cdots+n^2=\frac{n(n+1)(2n+1)}{6}$.
\item
שלב היסק: מתקיים
\begin{align*}
1^2+2^2+\cdots+(n+1)^2&=1^2+\cdots+n^2+(n+1)^2\\
&\underset{\text{הנחת האינדוקציה}}{=}\frac{n(n+1)(2n+1)}{6}+(n+1)^2\\
&=\frac{n(n+1)(2n+1)+6(n+1)^2}{6}\\
&=\frac{(n+1)\lsr{n(2n+1)+6(n+1)}}{6}\\
&=\frac{(n+1)(2n^2+7n+6)}{6}\\
&=\frac{(n+1)(n+2)(2n+3)}{6}\\
&=\frac{(n+1)\lsr{(n+1)+1}\lsr{2(n+1)+1}}{6},
\end{align*}
כנדרש.
\end{itemize}
\end{solbox}
\item
בארץ עוץ יש מטבעות בשווי של שלושה דולרי-עוץ וחמישה דולרי-עוץ.
הראו כי כל סכום שלם הגדול ממש משבעה דולרי-עוץ ניתן לפרוט
למטבעות של שלושה דולרי-עוץ וחמישה דולרי-עוץ.
\hint{נסחו את הטענה כהצגה $n=3a+5b$ עם $a,b$ שלמים אי-שליליים. בשלב המעבר החליפו כמה מטבעות בהצגה של $n$ כך שהסכום יגדל באחד, ופצלו למקרים לפי סוגי המטבעות שיש בה.}
\begin{solbox}
נתרגם את הבעיה לשפה מתמטית: כל מספר שלם
$n\geq8$
ניתן להצגה
$n=3a+5b$
כאשר
$a,b$
שלמים אי-שליליים.
\begin{itemize}
\item
שלב הבסיס: נציג
$8=3\cdot1+5\cdot1$,
כלומר, הטענה נכונה עבור
$n=8$.
\item
כעת נניח כי הטענה נכונה עבור איזשהו
$n\geq8$,
כלומר, נניח שקיימים שלמים אי-שליליים
$a,b$
כך ש-
$n=3a+5b$.
\item
שלב היסק: ע"פ הנחת האינדוקציה
$$.n+1=3a+5b+1$$
אם
$b\geq1$,
אז
$$,n+1=3a+5b+1=3a+5b+6-5=3(a+2)+5(b-1)$$
וזו הצגה רצויה. ואם
$b=0$,
אז, מכיוון ש-
$n\geq8$,
נסיק כי
$a\geq3$,
ולכן
$$,n+1=3a+1=3a+10-9=3(a-3)+5\cdot2$$
גם הפעם התקבלה הצגה רצויה.
\end{itemize}
\end{solbox}
\item
לכל טבעי
$n$:
$1+3+5+\cdots+(2n-1)=n^2$.
\begin{solbox}
נוכיח באינדוקציה כי לכל
$n$
טבעי סכום
$n$
המספרים האי-זוגיים הראשונים הוא
$n^2$.
\begin{itemize}
\item
בסיס האינדוקציה: עבור
$n=1$
הסכום הוא
$1=1^2$,
ולכן הטענה נכונה.
\item
שלב ההנחה: נניח כי
$1+3+\cdots+(2n-1)=n^2$
עבור איזשהו
$n$
טבעי.
\item
שלב היסק: האיבר הבא בסכום הוא
$2(n+1)-1=2n+1$,
ולכן
\begin{align*}
1+3+\cdots+(2n-1)+(2n+1)&\underset{\text{הנחת האינדוקציה}}{=}n^2+2n+1=(n+1)^2,
\end{align*}
כנדרש.
\end{itemize}
\end{solbox}
\end{enumerate}
